\chapter{Introduction}

Artificial Intelligence (AI) is a field that is rapidly evolving and has the potential to revolutionize various industries. Before delving into the technical aspects of AI, it is important to understand the concept of intelligence and how it can be measured.

\section{Intelligence Tests}

Scientists have tried to decide what differentiates AI from computer algorithms; i.e. what is the essence of intelligence? This question has been debated for decades, and there are different tests that have been proposed to determine whether a machine can be considered intelligent or not. The importance of these tests is shown in the Chinese Room thought experiment.
\begin{itemize}
    \item \ul{Chinese Room Experiment}:
    
    The Chinese Room thought experiment involves a person who does not understand Chinese sitting in a room with a set of instructions for how to respond to Chinese characters. The person receives Chinese characters as input and uses the instructions to produce appropriate responses, which are then sent back out of the room.
    
    From the outside, it may appear that the person in the room understands Chinese, but in reality, they are simply following a set of rules without any comprehension of the language. This thought experiment raises questions about whether a machine can truly understand language or if it is simply following a set of instructions.
\end{itemize}

Some of the most well-known tests include:
% // TODO: Memorizar Intelligence Tests
\begin{itemize}
    \item \ul{Turing Test}:
    
    In the Turing Test, a human evaluator interacts with both a machine and a human without knowing which is which. If the evaluator cannot reliably distinguish between the machine and the human, then the machine is said to have passed the Turing Test and is considered to be intelligent.

    It only tests for the ability to mimic human behavior, not for true understanding or consciousness.
    \item \ul{Winograd Test}:
    
    The test involves a series of sentences known as ``Winograd Schemas''. They include ambiguous pronouns that require common sense and contextual understanding to resolve.

    This test is not about factual knowledge but rather about understanding context.
    \item \ul{Lovelace Test}:
    
    The Lovelace Test is designed to evaluate a machine's ability to create something original and surprising.

    This test is extremely difficult for machines to pass, as they often just combine existing ideas rather than creating something truly new.
    \item \ul{Metzinger Test}:
    
    The Metzinger Test is designed to evaluate a machine's ability to have a subjective experience, a sense of ``self, I'', and consciousness. There is no definitive method for conducting the test in practice, since subjective experience cannot be observed from the outside.
\end{itemize}


\section{Initial Definitions}

From more general concepts to more specific ones, where each one builds upon the previous, some important terms in the field of AI include:
\begin{itemize}
    \item \ul{Data Science}: Field that turns data into knowledge and decisions.
    \item \ul{Artificial Intelligence}: Field that aims to create machines that can perform tasks that typically require human intelligence.
    \item \ul{Machine Learning}: Subfield of AI that focuses on training machines to learn from data and improve their performance over time without being explicitly programmed. 
    \item \ul{Neural Networks}: A type of (usually supervised) machine learning model inspired by the structure and function of the human brain. They consist of layers of interconnected nodes (neurons) that process and transmit information.
    \item \ul{Deep Learning}: Neural networks with multiple layers. Deep learning has been particularly successful in tasks such as image recognition, natural language processing, and speech recognition.
\end{itemize}